\documentclass[12pt]{article}
\usepackage{hyperref} 
\usepackage{booktabs}  
\usepackage[a4paper, margin=1in]{geometry}
\usepackage{titlesec}
\usepackage{enumitem}
\usepackage{setspace}
\usepackage{amsmath}
\usepackage{listings}
\usepackage{xcolor}

\setstretch{1.2}

\titleformat{\section}{\large\bfseries}{\thesection.}{1em}{}
\titleformat{\subsection}{\normalsize\bfseries}{\thesubsection}{1em}{}

% Listing style for OCaml
\lstset{
    language=Caml,
    basicstyle=\ttfamily\small,
    keywordstyle=\color{blue},
    commentstyle=\color{gray},
    stringstyle=\color{red},
    breaklines=true,
    frame=single,
    showstringspaces=false
}

\title{\textbf{Final Project Report:\\ Leveraging Meta-Programming in Functional Programming}}
\author{}
\date{}

\begin{document}

\maketitle

\noindent
\textbf{Project Title:} Leveraging Meta-Programming in Functional Programming \\
\textbf{Team Members:} Pratham Chawdhry (IMT2022068) and Vaibhav Mittal (IMT2022126) \\
\textbf{Language:} OCaml \\
\textbf{Build System:} Dune \\
\textbf{Github: } 
\href{https://github.com/Vaibhavvv24/PL_Project_MP/tree/main}{\textit{Github Repository}}.

\vspace{1em}

\section{Theory: Meta-Programming \& Functional Programming}

Meta-programming refers to the technique of writing programs that can generate, manipulate, or transform other programs (or themselves) at compile-time or runtime. Instead of solving a problem directly, meta-programs operate on code as data, enabling automation, abstraction, and optimization. Functional Programming (FP) provides a particularly robust foundation for this, as its features align naturally with the requirements of program transformation.

\subsection{Intersection of Meta-Programming and Functional Programming}
Functional programming languages treat computation as the evaluation of mathematical functions, avoiding mutable state. When combined with meta-programming, FP languages like OCaml excel because of their expressive type systems and syntactic capabilities.

Key functional characteristics that empower meta-programming include:
\begin{itemize}
    \item \textbf{Algebraic Data Types (ADTs):} Provide a rigorous and precise way to define Abstract Syntax Trees (ASTs) for representing code structures.
    \item \textbf{Pattern Matching:} Allows for concise, exhaustive destructuring of ASTs, making it significantly easier to write compiler passes, optimizers, and code generators.
    \item \textbf{First-Class Functions and Closures:} Enable high-order generation of program semantics and delayed evaluation.
    \item \textbf{Immutability:} Ensures safety when applying multiple transformations across an AST---each pass creates a new optimized tree rather than mutating state, preventing side-effect related compiler bugs.
\end{itemize}

\subsection{Types of Meta-Programming (Project Scope)}

There are two primary paradigms in meta-programming:
\begin{itemize}
    \item \textbf{Runtime Meta-Programming:} Programs evaluate, manipulate, and optimize their data structures or behavior dynamically while the host application is executing.
    \item \textbf{Compile-time Meta-Programming:} Code is generated or transformed before execution, such as OCaml's \textit{ppx} extensions. 
\end{itemize}

\subsection{Runtime Meta-Programming }
For this project, our primary focus is Runtime Meta-Programming. This involves treating program logic as data structures that can be dynamically inspected and optimized. In functional architectures, this commonly involves defining Domain-Specific Languages (DSLs) as deep embeddings using ADTs. Once the DSL captures an expression at runtime, the system can apply dynamic multi-pass optimizations---such as constant folding, algebraic simplification, and strength reduction. 



\section{Problem Description}

The core objective of this project is to explore how functional programming (FP) features---specifically algebraic data types, pattern matching, and first-class functions---enable the creation of robust meta-programming systems.

This project specifically demonstrates how programs can be treated as data and transformed across different stages of execution, enabling sophisticated program manipulation techniques.

\subsection*{Learning Objectives}
\begin{itemize}
    \item Understand the principles of \textbf{Runtime Meta-Programming} via deep AST embeddings and \textbf{Compile-Time Meta-Programming} via syntax extensions.
    \item Explore how logical expressions can be structurally manipulated for runtime optimization and staged generation.
    \item Implement an AST-based multi-pass optimizer demonstrating significant computational performance gains.
    \item Build advanced meta-programming pipelines including compile-time instrumentation (PPX) and higher-order memoization.
\end{itemize}

\subsection*{Intended Outcomes}
\begin{itemize}
    \item A high-performance, modular OCaml engine illustrating both runtime and compile-time meta-programming integrations.
    \item A structurally safe environment utilizing static typing and immutability for mathematically sound AST mutations.
    \item 
A specialized compiler phase designed to translate optimized DSL expressions directly into native, high-efficiency OCaml source code. By bypassing the interpretation layer entirely, this system bridges the gap between high-level symbolic logic and raw machine performance, allowing for "near-zero" runtime overhead.
    \item Practical implementation of \texttt{ppxlib} macros to demonstrate generative programming, using source-code instrumentation to automate "boilerplate" tasks during the compilation phase.
\end{itemize}

\section{Solution Outline}

Our final infrastructure showcases the complete meta-programming lifecycle:

\begin{enumerate}
    \item \textbf{Core FP Layer:}  
    A foundational utility layer leveraging first-class functions, lexical closures, and higher-order combinators (map, fold, filter) as the primary building blocks for the system’s execution engine.

    \item \textbf{Runtime Meta-System (DSL Engine):}
    \begin{itemize}
       
        \item \textbf{AST Representation:}  
A specialized tree structure for math expressions that prevents "broken" or invalid formulas. This ensures that the optimizer can safely rearrange and simplify calculations without risk of failure.

        \item \textbf{Multi-Pass Optimizer:}  
        A transformation pipeline that recursively applies optimization rules until achieving fixed-point convergence, ensuring maximal reduction of the expression tree.
        \begin{itemize}
            \item \textit{Constant Folding:} Pre-computation of deterministic sub-expressions.
            \item \textit{Algebraic Simplification:} Application of mathematical identity, zero, and absorption laws.
            \item \textit{Strength Reduction:} Replacement of computationally expensive operations (e.g., multiplication) with cheaper equivalents (e.g., recursive addition).
        \end{itemize}

        \item \textbf{Interpreter \& Evaluator:}  
An execution engine that processes math expressions by breaking them into smaller parts. It manages variable assignments in real-time and can pre-solve sub-calculations (Partial Evaluation) to improve efficiency.

        \item \textbf{Analysis \& Visualization:}
        A structural observability layer for debugging, providing pretty-printing capabilities and AST metrics (node count, tree depth) to track transformation impact across optimization passes.
    \end{itemize}

    \item \textbf{Advanced Meta-Programming Modules (Post-MidEval Stage):}
    \begin{itemize}
        \item \textbf{Offline Staged Compilation (CodeGen):}  
        A translation layer in \texttt{codegen.ml} that transforms optimized ASTs into native OCaml source code strings. It features precedence-aware parenthesization and automated free-variable detection to generate standalone native functions.
        
        \item \textbf{Compile-Time PPX Logging:}  
        A syntax rewriter implemented using \texttt{ppxlib} that automatically instruments functions with entry/exit logs (annotated via \texttt{[\%\%log let ...]}), demonstrating non-invasive instrumentation.
        
        \item \textbf{Higher-Order Memoization:}  
        A runtime meta-transformation that dynamically augments existing recursive functions with caching logic, showcasing the power of higher-order functions to optimize recursive patterns.
        
        \item \textbf{Meta-Transformation Taxonomy:}  
        A formal classification registry in \texttt{optimizer.ml} that categorizes every transform into Optimizing, Structural, or Instrumentation classes for structural observability.
    \end{itemize}

    \item \textbf{Benchmarking \& Validation:}
    \begin{itemize}
        \item \textbf{Performance Analytics:} High-precision scripts to quantify execution speedups and memory overhead reduction post-optimization.
        \item \textbf{Scalability Testing:} Comparative analysis of engine throughput when processing high-depth expression trees versus unoptimized recursive calls.
    \end{itemize}
\end{enumerate}

\section{Main References}

\begin{enumerate}
    \item \href{https://maestroinc.medium.com/meta-programming-in-python-28d5ace0223c}{\textit{Meta-Programming}}. \textbf{Relevance:} Provided conceptual grounding in runtime meta-programming techniques, particularly dynamic code transformation and reflection, which informed the design of our AST-based DSL transformation and optimization pipeline.
     \item \href{https://ocaml.org/docs}{\textit{ocaml documentation}}. \textbf{Relevance:} Provided conceptual grounding in ocaml , particularly higher order functions and pattern matching.
   
    \item \textit{ppxlib Tool Documentation} (OCaml Community).\\
    \href{https://ocaml-ppx.github.io/ppxlib/}{https://ocaml-ppx.github.io/ppxlib/}\\
    \textbf{Relevance:} The official tool literature detailing OCaml's Abstract Syntax Tree rewrite structures, which serves as the blueprint for writing our automated compile-time syntax extensions.
\end{enumerate}

\section{Final Progress and Accomplishments}

\subsection*{Final Milestones Completed}
\begin{itemize}
    \item Core functional programming abstraction layer implemented.
    \item DSL AST definition designed and integrated.
    \item Multi-pass optimizer pipeline fully functional with fixed-point convergence.
    \item \textbf{Offline Staged Compilation (CodeGen)} fully implemented and verified via \texttt{app\_codegen.exe}.
    \item \textbf{Compile-Time PPX} rewriter completed using \texttt{ppxlib} (\texttt{app\_logger.exe}).
    \item \textbf{Higher-Order Memoization} system integrated and benchmarked.
    \item Final unified CLI demo system developed (\texttt{demo/demo.ml}).
\end{itemize}

\subsection*{Key Achievements}
The project successfully bridges runtime and compile-time meta-programming. Through recursive AST manipulations, we achieved speedups exceeding \textbf{18,000$\times$} for arithmetic evaluations, while the CodeGen module allows for "Zero-Overhead" native execution. The system effectively demonstrates that OCaml is an ideal platform for building self-optimizing software architectures.

\subsection*{Challenges Overcome}
\begin{itemize}
    \item \textbf{Operator Precedence:} Solved by implementing a recursive parenthesizer in the CodeGen module to ensure generated OCaml code like \texttt{((score +. 5.) /. 2.)} remains valid.
    \item \textbf{Mutual Rule Triggering:} Prevented infinite loops in the optimizer by establishing a strict ordering for algebraic and constant-folding passes.
    \item \textbf{PPX Contexts:} Understanding the \texttt{ppxlib} context and location tracking to ensure generated code doesn't break source maps.
\end{itemize}


\section{Detailed Results and Performance Analysis}

The system was evaluated using the full suite of demonstrations and benchmarks. All results below were generated via the \texttt{make all} command.

\subsection{Functional Verification: Core FP Layer}
The \texttt{app\_core\_fp.ml} suite verified the foundational functional techniques:
\begin{itemize}
    \item \textbf{Closures:} \texttt{make\_adder 7} applied to 3 yielded \textbf{10}.
    \item \textbf{Function Pipelines:} A pipeline of \texttt{[+3, *2, +1]} applied to 4 correctly returned \textbf{15}.
    \item \textbf{Dynamic Generation:} The \texttt{dynamic\_power 3} (cube) function returned \textbf{125} for input 5.
    \item \textbf{HOFs:} \texttt{fold\_left (+)} over numbers 1--10 correctly calculated the sum \textbf{55}.
\end{itemize}

\subsection{Benchmark 1: Optimizer Performance}
Measured using a linear chain of 5,000 additions (10,001 AST nodes). The constant-folding optimizer reduced this entire tree to a single constant node.
\begin{itemize}
    \item \textbf{Unoptimized Execution:} 0.056702 s
    \item \textbf{Optimized Execution:} 0.000003 s
    \item \textbf{Final Speedup:} \textbf{18,294.2$\times$}
\end{itemize}

\subsection{Benchmark 2: Algebraic Simplification}
Evaluated a redundant symbolic expression: $(((x + 0) \times 1) + (0 \times (y + 99)))$.
\begin{itemize}
    \item \textbf{Unoptimized:} 0.001436 s
    \item \textbf{Optimized:} 0.000312 s
    \item \textbf{Final Speedup:} \textbf{4.6$\times$}
\end{itemize}

\subsection{Benchmark 3: Higher-Order Memoization}
Applied meta-caching to a recursive \texttt{fib(35)} call.
\begin{itemize}
    \item \textbf{Naive Execution:} 0.049968 s (\textbf{29,860,703 calls})
    \item \textbf{Memoized Execution:} 0.000002 s (\textbf{69 calls})
    \item \textbf{Final Speedup:} \textbf{26,197.6$\times$}
    \item \textbf{Call Reduction:} 99.99\% of naive calls eliminated.
\end{itemize}

\subsection{Step-by-Step Optimization Pipeline}
The multi-pass pipeline was verified to correctly sequence transformations:
\begin{enumerate}
    \item \textbf{Original:} \texttt{((y * 0) + (5 + 10))} (7 nodes)
    \item \textbf{After Fold:} \texttt{((y * 0) + 15)} (5 nodes)
    \item \textbf{After Simplify:} \texttt{15} (1 node)
\end{enumerate}

\subsection{Partial Evaluation}
The system demonstrated symbolic reduction on a complex polynomial:
\begin{itemize}
    \item \textbf{Symbolic:} $((1 + (2 \times x)) + (3 \times (x \times x)))$
    \item \textbf{Partial Eval ($x=0$):} Reduced to \texttt{1}
    \item \textbf{Full Eval ($x=2$):} Correctly calculated \textbf{17}.
\end{itemize}

\section{Work Distribution (Final)}

\subsection*{Member 1 (IMT2022126)}
\begin{itemize}
    \item Developed core FP foundations and DSL environment bindings (\texttt{src/core\_fp/}, \texttt{src/runtime\_meta/dsl.ml}).
    \item Engineered the \textbf{Offline Staged Compilation (CodeGen)} system (\texttt{src/runtime\_meta/codegen.ml}).
    \item Built the final integrated system demo and CLI walkthrough (\texttt{demo/demo.ml}).
    \item Implemented Core FP demonstrations (\texttt{examples/app\_core\_fp.ml}).
\end{itemize}

\subsection*{Member 2 (IMT2022068)}
\begin{itemize}
    \item Developed the multi-pass AST Transformation and Optimization engine (\texttt{src/runtime\_meta/optimizer.ml}).
    \item Engineered the \textbf{Compile-Time PPX Logger} extension (\texttt{src/compiletime\_meta/}).
    \item Developed the \textbf{Higher-Order Memoization} and Benchmark suites (\texttt{benchmarks/}).
    \item Implemented the AST Visualization and Meta-Transformation Taxonomy registry.
\end{itemize}

\vspace{1em}

\section*{Appendix (Technical Details)}

\subsection*{Offline Staged Compilation Output}
Test Case 1: Complex Let-Binding \& Strength Reduction.
\begin{lstlisting}[language=Caml]
(* Original: (let z = (x * 2) in ((z + y) * ((10 * 0) + (15 - 5)))) *)
let compute_value x y =
  (let z = (x + x) in
  ((z + y) * 10))
\end{lstlisting}

Test Case 4: Floating Point Generation.
\begin{lstlisting}[language=Caml]
let compute_average score =
  ((score +. 5.) /. 2.)
\end{lstlisting}

\subsection*{Meta-Transformation Taxonomy Registry}
\begin{tabular}{lll}
    \toprule
    \textbf{Transform Name} & \textbf{Class} & \textbf{Description} \\
    \midrule
    \texttt{constant\_folding} & Optimization & Pre-computes constants \\
    \texttt{algebraic\_simplification} & Optimization & Identity/Zero laws \\
    \texttt{strength\_reduction} & Structural & x*2 $\to$ x+x \\
    \texttt{full\_optimizer} & Optimization & Fixed-point convergence \\
    \bottomrule
\end{tabular}

\section*{Conclusion}
The final phase of this project successfully bridged the gap between runtime interpretation and native performance. By combining functional power with meta-programming techniques, we demonstrated a system that is both highly abstract and extremely performant, achieving speedups of up to \textbf{26,000$\times$} through symbolic manipulation.

\end{document}
